\section*{الفصل \ref{ch:g12:catalysis} --- الحفز}

\begin{solution}{exo:g12:catalysis:1}
\omterm{def:g7:chemical-reactions:reactant}{المتفاعل}: \ce{H2O2}. \omterm{def:g7:chemical-reactions:reactant}{النواتج}: \ce{H2O} و\ce{O2}. \omterm{def:g12:catalysis:catalyst}{الحفّاز}: \ce{I-}
(مستهلك في الخطوة 1، ومُعاد في الخطوة 2). الوسيط: \ce{IO-} (يتكوّن في
الخطوة 1، ويُستهلك في الخطوة 2).
\end{solution}

\begin{solution}{exo:g12:catalysis:2}
متجانس؛ غير متجانس؛ غير متجانس؛ متجانس.
\end{solution}

\begin{solution}{exo:g12:catalysis:3}
الطريق غير المحفَّز. والمستويان الابتدائي والنهائي هما نفسهما: فتتكوّن \omterm{def:g7:chemical-reactions:reactant}{النواتج}
نفسها، مع إطلاق الطاقة نفسها؛ ولا يختلف إلا الطريق.
\end{solution}

\begin{solution}{exo:g12:catalysis:4}
يُستهلك في خطوة ويُعاد في أخرى: فهو يظهر في
طرفي مجموع الخطوات، فيُحذف منهما.
\end{solution}

\begin{solution}{exo:g12:catalysis:5}
يجري التفاعل على السطح: فكلما كبر السطح، كثرت
\omterm{def:g7:atoms-and-molecules:molecule}{الجزيئات} التي يمكنها أن تتفاعل في آن واحد، للكتلة نفسها من \omterm{def:g12:catalysis:catalyst}{الحفّاز}
(الغالي في الغالب).
\end{solution}

\begin{solution}{exo:g12:catalysis:6}
المجموع: \ce{2Fe^{3+} + 2H2O2 + 2Fe^{2+} + 2H+ -> 2Fe^{2+} + O2 + 2H+ +
2Fe^{3+} + 2H2O}. \omterm{def:g9:ions:ion}{وأيونات} الحديد \omterm{def:g9:ions:ion}{وأيونات} الهيدروجين تظهر في الطرفين
فتُحذف: \ce{2H2O2 -> 2H2O + O2}.
\end{solution}

\begin{solution}{exo:g12:catalysis:7}
نصف \qty{72}{mL} يساوي \qty{36}{mL}: ويُبلغ بعد نحو \qty{35}{min}
دون حفّاز و\qty{3.5}{min} معه. \omterm{def:g12:catalysis:catalyst}{فالحفّاز} يقصّر المدة عشر
مرات.
\end{solution}

\begin{solution}{exo:g12:catalysis:8}
يجري التفاعل على سطح البلاتين؛ والرصاص يغطيه
ويبقى، فلا تعود الغازات تبلغ \omterm{def:g9:periodic-table-first-look:metal}{الفلز}. \omterm{def:g12:catalysis:catalyst}{فالحفّاز}
«مسموم»: موجود لكنه بلا فائدة.
\end{solution}

\begin{solution}{exo:g12:catalysis:9}
حتى نحو \qty{37}{\celsius}، يتسارع التفاعل بارتفاع
درجة الحرارة، ككل تفاعل. وفوق ذلك، يفقد \omterm{def:g12:catalysis:enzyme}{الإنزيم}، وهو بروتين،
شكله ومعه نشاطه؛ وعند \qty{80}{\celsius} يتلف.
\end{solution}

\begin{solution}{exo:g12:catalysis:10}
\ce{2CO + 2NO -> 2CO2 + N2}: C 2 = 2، O 4 = 4، N 2 = 2. يتأكسد أحادي أكسيد
الكربون بأحادي أكسيد النيتروجين، الذي يُختزل: فكل ملوِّث يزيل
الآخر.
\end{solution}

\begin{solution}{exo:g12:catalysis:11}
$n(\ce{O2}) = 0.072 / 24.1 = \qty{2.99e-3}{mol}$؛
$n(\ce{H2O2}) = 2 \times \num{2.99e-3} = \qty{5.98e-3}{mol}$ في
\qty{10.0}{mL}: $c = \qty{0.598}{mol/L}$.
\end{solution}

\begin{solution}{exo:g12:catalysis:12}
\ce{C2H5OH -> C2H4 + H2O}، وهو حذف؛ \ce{C2H5OH -> CH3CHO + H2}.
\omterm{def:g12:catalysis:selective}{والحفّاز الانتقائي} يسرّع تفاعلًا واحدًا من عدة تفاعلات ممكنة أكثر بكثير
من غيره، فيحدد بذلك \omterm{def:g7:chemical-reactions:reactant}{الناتج}.
\end{solution}

\begin{solution}{exo:g12:catalysis:13}
لا يغيّر \omterm{def:g12:catalysis:catalyst}{الحفّاز} \omterm{def:g10:reaction-progress-table:system}{الحالة النهائية}: فالمنحنيان مع \omterm{def:g12:catalysis:catalyst}{الحفّاز}
ودونه ينتهيان عند الحجم نفسه من الأكسجين. فيُحصل على الكمية نفسها من
\omterm{def:g7:chemical-reactions:reactant}{الناتج}، لكن أسرع فحسب.
\end{solution}

\begin{solution}{exo:g12:catalysis:14}
$\num{6.02e23} / \num{5e6} = \qty{1.2e17}{s}$، أي نحو أربعة مليارات سنة.
ولإنجاز ذلك في ثانية واحدة: \num{1.2e17} \omterm{def:g7:atoms-and-molecules:molecule}{جزيء} من \omterm{def:g12:catalysis:enzyme}{الإنزيم}، أي
\qty{2e-7}{mol}.
\end{solution}

\begin{solution}{exo:g12:catalysis:15}
الأمونيا \ce{NH3}: $150 \times 17.0 / 14.0 = \qty{182}{Mt}$. \omterm{def:g7:atoms-and-molecules:atom}{وذرات}
النيتروجين: $\num{1.50e14} / 14.0 = \qty{1.07e13}{mol}$، أي
\qty{5.36e12}{mol} من \ce{N2}.
\end{solution}

\begin{solution}{pb:g12:catalysis:1}
\textbf{1.} \ce{2H2O2 -> 2H2O + O2}.

\textbf{2.} نعم: فبيروكسيد الهيدروجين هو \omterm{def:g11:redox:oxidant}{مؤكسد} المزدوجة
$\ce{H2O2}/\ce{H2O}$ \omterm{def:g11:redox:oxidant}{ومختزل} المزدوجة $\ce{O2}/\ce{H2O2}$؛ فهو
يؤكسد نفسه ويختزلها.

\textbf{3.} التفاعل دون حفّاز بطيء جدًا.

\textbf{4.} $20 / 22.4 = \qty{0.893}{mol}$.

\textbf{5.} مولان من بيروكسيد الهيدروجين لكل \omterm{def:g10:the-mole:amount}{مول} من الأكسجين:
\qty{1.79}{mol}.

\textbf{6.} \qty{1.79}{mol/L}.

\textbf{7.} $M = 2 \times 1.0 + 2 \times 16.0 = \qty{34.0}{g/mol}$;
$C_m = 1.79 \times 34.0 = \qty{60.7}{g/L}$.

\textbf{8.} النصف: \qty{0.893}{mol/L}.

\textbf{9.} الكاتالاز: \omterm{def:g12:catalysis:enzyme}{إنزيم}، \omterm{def:g6:solutions-and-solubility:solute}{مذاب}، أي حفّاز حيوي؛
\omterm{def:g9:ions:ion}{وأيونات} الحديد (III): \omterm{def:g12:catalysis:homogeneous}{حفز متجانس}.

\textbf{10.} \ce{2Fe^{3+} + H2O2 -> 2Fe^{2+} + O2 + 2H+}
و\ce{2Fe^{2+} + H2O2 + 2H+ -> 2Fe^{3+} + 2H2O}؛ ومجموعهما
\ce{2H2O2 -> 2H2O + O2}.

\textbf{11.} مع \omterm{def:g12:catalysis:catalyst}{الحفّاز} يرتفع الحجم أسرع بكثير؛ ويستقر المنحنيان
عند الحجم النهائي نفسه.

\textbf{12.} تتحول \omterm{def:g9:ions:ion}{أيونات} الحديد (III) (البرتقالية) إلى حديد (II) (أخضر
باهت) في الخطوة الأولى وتعود إلى حديد (III) في الثانية؛ وفي
النهاية يعود الحديد كله حديدًا (III).

\textbf{13.} أتلف الغلي \omterm{def:g12:catalysis:enzyme}{الإنزيم}.

\textbf{14.} $0.100 \times 20 = \qty{2.0}{L}$.

\textbf{15.} $n(\ce{O2}) = 0.100 \times 0.893 = \qty{0.0893}{mol}$، أي
$0.0893 \times 24.1 = \qty{2.15}{L}$.

\textbf{16.} التفكك البطيء يطلق الأكسجين، الذي يرفع
الضغط في قارورة محكمة الإغلاق.

\textbf{17.} $1.50 / 2 = \qty{0.750}{mol}$ من الأكسجين في اللتر، أي
$0.750 \times 22.4 = \qty{16.8}{L}$: «16.8 حجمًا».

\textbf{18.} $(1.79 - 1.50)/1.79 \approx \qty{16}{\percent}$.

\textbf{19.} \qty{1.79}{mol/L}.
\end{solution}
